% BHU PYQ -- Manually transcribed & error-corrected
\documentclass[11pt,a4paper]{article}

\usepackage[a4paper,top=2.5cm,bottom=2.5cm,left=2.8cm,right=2.8cm,headheight=14pt]{geometry}
\usepackage{amsmath,amssymb}
\usepackage[T1]{fontenc}
\usepackage[utf8]{inputenc}
\usepackage{lmodern}
\usepackage{microtype}
\usepackage{fancyhdr}
\usepackage{enumitem}
\usepackage{hyperref}
\usepackage{lastpage}
\usepackage{parskip}

\hypersetup{colorlinks=true,urlcolor=black,linkcolor=black,
  pdftitle={CHB-401: Inorganic & Organic Chemistry-II -- BHU PYQ},pdfauthor={Banaras Hindu University}}

\pagestyle{fancy}\fancyhf{}
\renewcommand{\headrulewidth}{0.4pt}\renewcommand{\footrulewidth}{0.4pt}
\fancyhead[L]{\small\textit{Banaras Hindu University}}
\fancyhead[R]{\small\textit{CHB-401: Inorganic \& Organic Chemistry-II}}
\fancyfoot[C]{\small Page \thepage\ of \pageref{LastPage}}

\newlist{parts}{enumerate}{2}
\setlist[parts,1]{label=(\alph*),leftmargin=*,itemsep=5pt,topsep=3pt}
\setlist[parts,2]{label=(\roman*),leftmargin=*,itemsep=3pt,topsep=2pt}
\newcommand{\pts}[1]{\hfill\textit{[#1]}}

\begin{document}

\begin{center}
  {\large\bfseries BANARAS HINDU UNIVERSITY}\\[2pt]
  {\normalsize B.Sc. (Hons.) Semester IV Examination, 2022-23}\\[6pt]
  \rule{0.8\linewidth}{0.5pt}\\[6pt]
  {\large\bfseries Paper No. CHB-401: Inorganic \& Organic Chemistry-II}\\[4pt]
  {\normalsize Time: 3 Hours\hfill Maximum Marks: 70}
\end{center}
\rule{\linewidth}{0.4pt}
\small
\noindent\textit{Instructions: Answer five questions in total, including Question~1 from each section which is compulsory. Use separate answer books for Section~A and Section~B.}
\normalsize
\bigskip

\begin{center}
  {\large\bfseries SECTION A}\\[2pt]
  {\normalsize (Inorganic Chemistry-II -- Marks: 35)}
\end{center}
\rule{0.4\linewidth}{0.2pt}
\smallskip

\noindent\textbf{Question 1.} Answer all parts of the following: \pts{5 $\times$ 2 = 10}
\begin{parts}
  \item The water solubility of lithium halides is in the order $\text{LiBr} > \text{LiCl} > \text{LiF}$. However, this trend is reversed in the case of analogous silver halides. Explain.
  \item What are the differences in characteristic properties of the blue colored solution and bronze colored phase obtained by addition of alkali metals in liquid ammonia?
  \item What are the limitations of the Effective Atomic Number (EAN) rule?
  \item Why do the $4\text{d}$ and $5\text{d}$ transition series elements show very similar ionic radii for corresponding elements? Explain.
  \item Predict the relative coordination tendencies of $\text{F}^-$ and $\text{I}^-$ towards $\text{Fe}^{3+}$ and $\text{Hg}^{2+}$ using HSAB principles.
\end{parts}

\medskip\rule{\linewidth}{0.2pt}

\noindent\textbf{Question 2.}
\begin{parts}
  \item Give two examples each of protic and aprotic solvents and write autoionization reactions for them. \pts{4}
  \item Predict the product when $\text{(CH}_3)_3\text{N-PF}_2$ is reacted with: (i) $\text{BH}_3$ and (ii) $\text{BF}_3$. Explain the reactivity difference between the two cases. \pts{4}
  \item State the physical properties of a solvent that determine its utility in chemical reactions. \pts{4}
\end{parts}

\medskip\rule{\linewidth}{0.2pt}

\noindent\textbf{Question 3.}
\begin{parts}
  \item Write one example each of the following types of octahedral complexes. Draw the various stereoisomers and specify the enantiomeric pairs in each case: \pts{4}
  \begin{parts}
    \item $[\text{Ma}_2\text{b}_2\text{c}_2]$
    \item $[\text{M(AA)}_2\text{b}_2]$
  \end{parts}
  \item Write the IUPAC names of the following coordination compounds: \pts{4}
  \begin{parts}
    \item $\text{K}_2[\text{Zn(NCS)}_4]$
    \item $[\text{Mn(acac)}_3]$
    \item $[(\text{NH}_3)_5\text{Co-NH}_2\text{-Co(NH}_3)_5]\text{Cl}_5$
  \end{parts}
  \item Why is $[\text{Ni(en)}_3]^{2+}$ more stable than $[\text{Ni(NH}_3)_6]^{2+}$ even though both are octahedral and nitrogen-donor complexes? Draw their structures. \pts{4}
\end{parts}

\medskip\rule{\linewidth}{0.2pt}

\noindent\textbf{Question 4.}
\begin{parts}
  \item Using Valence Bond Theory, predict the structure, hybridization, and spin-only magnetic moments of: $\text{Ni(CO)}_4$, $[\text{NiCl}_4]^{2-}$, $[\text{Ni(CN)}_4]^{2-}$, and $[\text{Ni(NH}_3)_6]^{2+}$. \pts{6}
  \item Write the electronic configuration and number of unpaired electrons in the ground state of the following lanthanide ions: $\text{Ce}^{3+}$, $\text{Gd}^{3+}$, $\text{Eu}^{3+}$, and $\text{Tb}^{3+}$. \pts{3}
  \item Predict the feasibility of the following reactions in water and liquid ammonia solvents along with reasons: \pts{3}
  \begin{parts}
    \item $\text{BaCl}_2 + 2\text{AgNO}_3 \to 2\text{AgCl} + \text{Ba(NO}_3)_2$
    \item $\text{AgCl} + \text{KNO}_3 \to \text{KCl} + \text{AgNO}_3$
  \end{parts}
\end{parts}

\newpage
\begin{center}
  {\large\bfseries SECTION B}\\[2pt]
  {\normalsize (Organic Chemistry-III -- Marks: 35)}
\end{center}
\rule{0.4\linewidth}{0.2pt}
\smallskip

\noindent\textbf{Question 5.} Answer all parts of the following: \pts{5 $\times$ 2 = 10}
\begin{parts}
  \item Write the chemical structures of $\alpha$- and $\beta$-D-(+)-glucopyranose.
  \item Discuss mutarotation of glucose in brief.
  \item How do epimers and anomers differ? Explain with examples.
  \item Give one method for the synthesis of Indigo dye from anthranilic acid.
  \item Why do malachite green and fluorescein develop color in solution? Explain.
\end{parts}

\medskip\rule{\linewidth}{0.2pt}

\noindent\textbf{Question 6.}
\begin{parts}
  \item Discuss the orientation of electrophilic aromatic substitution reactions in pyridine. Explain why substitution occurs preferentially at position 3. \pts{4}
  \item Explain why pyridine $N$-oxide is more reactive than pyridine towards both electrophilic and nucleophilic substitution reactions. \pts{4}
  \item Write the products in the following reactions: \pts{4}
  \begin{parts}
    \item $\text{Quinoline} \xrightarrow{\text{NaNH}_2, 100^\circ\text{C}} ?$
    \item $\text{Pyrrole} \xrightarrow{\text{CH}_3\text{I, KOH}} ?$
  \end{parts}
\end{parts}

\medskip\rule{\linewidth}{0.2pt}

\noindent\textbf{Question 7.}
\begin{parts}
  \item Discuss the synthesis of the following heterocyclic compounds with mechanisms: \pts{6}
  \begin{parts}
    \item Skraup quinoline synthesis
    \item Bischler-Napieralski isoquinoline synthesis
  \end{parts}
  \item Explain why anthracene is very reactive at the 9 and 10 positions. \pts{3}
  \item How will you establish the structure of indigotin? Explain its chemical degradation. \pts{3}
\end{parts}

\medskip\rule{\linewidth}{0.2pt}

\noindent\textbf{Question 8.}
\begin{parts}
  \item How will you prepare the following compounds from D-(+)-glucose? Write stepwise mechanisms: \pts{6}
  \begin{parts}
    \item Fructose (Osazone method)
    \item Mannose (Epimerization)
  \end{parts}
  \item Complete the following reactions: \pts{6}
  \begin{parts}
    \item $\text{D-Erythrose} + \text{CH}_3\text{NO}_2 \xrightarrow{\text{i. CH}_3\text{ONa, ii. H}_2\text{SO}_4} ?$
    \item $\text{D-Glucose} \xrightarrow{\text{i. NH}_2\text{OH, ii. Ac}_2\text{O, NaOAc, iii. NaOMe}} ?$
  \end{parts}
\end{parts}

\vfill
\rule{\linewidth}{0.4pt}
\begin{center}\small
  \href{https://bhu-exam.vercel.app/}{Click here to visit our website}\\[4pt]
  \href{https://me-aryan.vercel.app/}{Click here to visit our contributor}\\[4pt]
  \href{https://quashan-me.vercel.app/}{Click here to visit our contributor}
\end{center}

\end{document}
